\documentclass[12pt,a4paper]{article}

\usepackage[utf8]{inputenc}
\usepackage[T1]{fontenc}
\usepackage[brazil]{babel}
\usepackage{amsmath,amssymb}
\usepackage{algorithm}
\usepackage{algpseudocode}
\usepackage{geometry}

\geometry{margin=2.5cm}

\title{Verificação de Bipartição por Busca em Largura}
\author{Prof. Me. Luis Vinicius Costa Silva}
\date{}

\begin{document}

\maketitle

\section{Problema}

Considere um grafo simples e não direcionado

$$
G=(V,E),
$$

em que cada vértice representa um elemento e cada aresta representa
um conflito entre dois elementos.

Deseja-se determinar se os vértices podem ser divididos em dois grupos
de modo que nenhuma aresta conecte dois vértices pertencentes ao mesmo
grupo.

Esse problema é equivalente a determinar se o grafo é
\emph{bipartido}.

Um grafo é bipartido se seu conjunto de vértices pode ser particionado
em dois conjuntos disjuntos

$$
V=V_0\cup V_1,
\qquad
V_0\cap V_1=\varnothing,
$$

tal que toda aresta possui uma extremidade em $V_0$ e outra em $V_1$.

\section{Ideia da solução}

A verificação pode ser realizada por uma busca em largura
(\emph{Breadth-First Search}, BFS) combinada com uma
\emph{2-coloração} dos vértices.

Para cada vértice $v$, define-se uma variável de cor:

$$
cor[v]\in\{-1,0,1\},
$$

onde:

\begin{itemize}
\item $cor[v]=-1$: vértice ainda não colorido;
\item $cor[v]=0$: vértice pertencente ao grupo $0$;
\item $cor[v]=1$: vértice pertencente ao grupo $1$.
\end{itemize}

Quando a BFS percorre uma aresta $(u,v)$, devem ser satisfeitas as
seguintes condições:

\begin{enumerate}
\item se $v$ ainda não possui cor, atribui-se a ele a cor oposta
à de $u$;
\item se $v$ já possui cor e $cor[v]=cor[u]$, existe um conflito
e o grafo não é bipartido;
\item se $cor[v]\neq cor[u]$, a aresta é compatível com a
bipartição.
\end{enumerate}

A regra de atribuição pode ser escrita como

$$
cor[v]=1-cor[u].
$$

\section{Componentes desconectados}

É necessário executar a BFS a partir de todos os vértices que ainda
não foram coloridos.

Isso ocorre porque o grafo pode possuir mais de uma componente conexa.

O algoritmo segue a estrutura:

\begin{algorithm}
\caption{BFS para verificação de bipartição}
\begin{algorithmic}[1]

\Require Grafo $G=(V,E)$
\Ensure O grafo é bipartido ou não

\ForAll{$v\in V$}
\State $cor[v]\gets -1$
\EndFor

\ForAll{$s\in V$}
\If{$cor[s]\neq -1$}
\State \textbf{continue}
\EndIf


\State $cor[s]\gets 0$
\State Inserir $s$ na fila

\While{a fila não estiver vazia}
    \State $u\gets$ remover elemento da fila

    \ForAll{$v$ adjacente a $u$}
        \If{$cor[v]=-1$}
            \State $cor[v]\gets 1-cor[u]$
            \State Inserir $v$ na fila
        \ElsIf{$cor[v]=cor[u]$}
            \State \Return \textbf{DIVISÃO IMPOSSÍVEL}
        \EndIf
    \EndFor
\EndWhile


\EndFor

\State \Return \textbf{DIVISÃO POSSÍVEL}

\end{algorithmic}
\end{algorithm}

\section{Representação por matriz de adjacência}

Na primeira abordagem, o grafo é representado por uma matriz

$$
A\in\{0,1\}^{n\times n},
$$

definida por

$$
A_{ij}=
\begin{cases}
1,&\text{se }(i,j)\in E,\\
0,&\text{caso contrário}.
\end{cases}
$$

Como o grafo é não direcionado,

$$
A_{ij}=A_{ji}.
$$

Para verificar os vizinhos de um vértice $u$, percorre-se toda a linha
correspondente a $u$:

$$
A_{u0},A_{u1},\ldots,A_{u,n-1}.
$$

Um vértice $v$ é vizinho de $u$ quando

$$
A_{uv}=1.
$$

A parte principal da BFS pode ser representada pelo seguinte
pseudocódigo:

\begin{algorithm}
\caption{BFS utilizando matriz de adjacência}
\begin{algorithmic}[1]

\State $u\gets$ remover elemento da fila

\For{$v\gets0$ até $n-1$}


\If{$A[u][v]=0$}
    \State \textbf{continue}
\EndIf

\If{$cor[v]=-1$}
    \State $cor[v]\gets1-cor[u]$
    \State Inserir $v$ na fila

\ElsIf{$cor[v]=cor[u]$}
    \State \Return \textbf{DIVISÃO IMPOSSÍVEL}
\EndIf


\EndFor

\end{algorithmic}
\end{algorithm}

\subsection{Implementação em C}

A estrutura básica do grafo é:

\begin{verbatim}
#define MAX 100

typedef struct {
int n;
int adj[MAX][MAX];
} Grafo;
\end{verbatim}

A inserção de uma aresta não direcionada é:

\begin{verbatim}
void adicionar_aresta(Grafo *g, int u, int v) {
g->adj[u][v] = 1;
g->adj[v][u] = 1;
}
\end{verbatim}

A busca pelos vizinhos é realizada por:

\begin{verbatim}
for (int v = 0; v < g->n; v++) {


if (g->adj[u][v] == 0)
    continue;

if (cor[v] == -1) {
    cor[v] = 1 - cor[u];
    enfileirar(&fila, v);
}
else if (cor[v] == cor[u]) {
    return 0;
}


}
\end{verbatim}

\section{Representação por TAD de grafo com ponteiros}

Uma segunda abordagem consiste em representar o grafo por listas de
adjacência.

Cada vértice possui uma lista contendo seus vizinhos.

Um possível TAD é:

\begin{verbatim}
typedef struct No {
int vertice;
struct No *prox;
} No;

typedef struct {
No *inicio;
} Lista;

typedef struct {
int n;
Lista adj[MAX];
} Grafo;
\end{verbatim}

Nesse caso, para cada vértice $u$, a estrutura

$$
g->adj[u].inicio
$$

aponta para o primeiro elemento da lista de seus vizinhos.

Os demais vizinhos são acessados por meio do campo

$$
prox.
$$

\subsection{Percorrendo os vizinhos}

Na matriz de adjacência, os vizinhos são encontrados percorrendo uma
linha inteira da matriz.

Na lista de adjacência, percorre-se somente os elementos que realmente
representam arestas:

\begin{verbatim}
No *p = g->adj[u].inicio;

while (p != NULL) {


int v = p->vertice;

if (cor[v] == -1) {

    cor[v] = 1 - cor[u];
    enfileirar(&fila, v);

}
else if (cor[v] == cor[u]) {

    return 0;
}

p = p->prox;


}
\end{verbatim}

A lógica da BFS permanece exatamente a mesma.

A única diferença está na maneira como os vizinhos de $u$ são
obtidos.

\section{TAD da fila}

Independentemente da representação utilizada para o grafo, a BFS
necessita de uma fila.

Uma implementação simples pode ser definida por:

\begin{verbatim}
typedef struct {
int dados[MAX];
int inicio;
int fim;
} Fila;
\end{verbatim}

As operações fundamentais são:

\begin{itemize}
\item inicializar a fila;
\item verificar se a fila está vazia;
\item inserir um vértice;
\item remover um vértice.
\end{itemize}

A fila garante que os vértices sejam processados na ordem em que são
descobertos, característica fundamental da BFS.

\section{Exemplo}

Considere o grafo

$$
0-1-2-3.
$$

Sua matriz de adjacência é

$$
A=
\begin{pmatrix}
0&1&0&0\\
1&0&1&0\\
0&1&0&1\\
0&0&1&0
\end{pmatrix}.
$$

Inicialmente,

$$
cor=[-1,-1,-1,-1].
$$

A BFS começa pelo vértice $0$:

$$
cor[0]=0.
$$

Como $1$ é vizinho de $0$,

$$
cor[1]=1.
$$

Em seguida, $2$ recebe a cor oposta de $1$:

$$
cor[2]=0.
$$

Finalmente,

$$
cor[3]=1.
$$

Obtém-se:

$$
V_0=\{0,2\},
\qquad
V_1=\{1,3\}.
$$

Não existe aresta entre dois vértices do mesmo grupo. Portanto,

$$
\boxed{\text{DIVISÃO POSSÍVEL}}
$$

\section{Detecção de conflito}

Considere agora um triângulo:

$$
0-1-2-0.
$$

A BFS pode produzir:

$$
cor[0]=0,
\qquad
cor[1]=1,
\qquad
cor[2]=1.
$$

Ao analisar a aresta $(2,0)$, verifica-se:

$$
cor[2]=1
$$

e

$$
cor[0]=0.
$$

Nesse caso, não há conflito.

Entretanto, a aresta $(1,2)$ exige que os dois vértices tenham cores
diferentes. Como ambos possuem cor $1$, ocorre:

$$
cor[1]=cor[2].
$$

O algoritmo retorna imediatamente:

$$
\boxed{\text{DIVISÃO IMPOSSÍVEL}}
$$

Esse conflito corresponde à existência de um ciclo ímpar.

\section{Complexidade}

A complexidade depende da representação utilizada.

\subsection{Matriz de adjacência}

Para cada vértice processado, é necessário percorrer toda a linha da
matriz para encontrar seus vizinhos.

Assim, a complexidade é

$$
O(|V|^2).
$$

O espaço necessário para armazenar a matriz é

$$
O(|V|^2).
$$

Essa representação possui como vantagem o acesso direto:

$$
A_{uv}
$$

pode ser consultado em tempo

$$
O(1).
$$

\subsection{Lista de adjacência}

Na representação por listas, a BFS percorre apenas as arestas
existentes.

Consequentemente, a complexidade é

$$
O(|V|+|E|).
$$

O espaço necessário para armazenar o grafo é

$$
O(|V|+|E|).
$$

Essa representação é especialmente vantajosa para grafos esparsos,
nos quais

$$
|E|\ll |V|^2.
$$


\section{Comparação das representações}

As duas representações implementam exatamente o mesmo algoritmo de
BFS e utilizam a mesma regra de 2-coloração.

A diferença está na estrutura utilizada para encontrar os vizinhos.

\begin{center}
\begin{tabular}{lcc}
\hline
\textbf{Característica} &
\textbf{Matriz} &
\textbf{Lista de adjacência} \\
\hline
Armazenamento &
$O(V^2)$ &
$O(V+E)$ \\

Consulta de uma aresta &
$O(1)$ &
$O(\deg(u))$ \\

Obtenção dos vizinhos &
$O(V)$ &
$O(\deg(u))$ \\

Complexidade da BFS &
$O(V^2)$ &
$O(V+E)$ \\

Implementação &
Simples e direta &
Ponteiros e listas encadeadas \\

Mais adequada para &
Grafos densos &
Grafos esparsos \\
\hline
\end{tabular}
\end{center}

\subsection{Diferença na obtenção dos vizinhos}

Do ponto de vista algorítmico, não há diferença na estratégia de
resolução:

\[
\boxed{\text{BFS}+\text{2-coloração}}
\]

é a mesma nas duas representações.

O que muda é a operação de obtenção dos vizinhos de um vértice $u$.

Na matriz de adjacência, percorremos toda a linha correspondente a
$u$:

\[
\text{vizinhos}(u)
=
\{v\mid A_{uv}=1\}.
\]

Na lista de adjacência, percorremos somente os elementos armazenados
na lista de $u$:

\[
\text{vizinhos}(u)
=
\text{lista de adjacência de }u.
\]

Assim, em um grafo com poucos conflitos, a lista de adjacência evita
percorrer posições da matriz que correspondem a pares de vértices sem
arestas.

\subsection{Resumo}

A escolha pode ser resumida da seguinte forma:

\[
\boxed{
\begin{array}{ll}
\text{Matriz:} & O(V^2)\\[2mm]
\text{Lista:} & O(V+E)
\end{array}
}
\]

Portanto, para um grafo esparso, a lista de adjacência tende a ser
mais eficiente. Para grafos pequenos ou quando se deseja uma
implementação mais simples, a matriz de adjacência pode ser uma
alternativa conveniente.

\section{Conclusão}

A verificação de bipartição pode ser realizada de maneira eficiente
por BFS e 2-coloração.

O algoritmo atribui uma das duas cores a cada vértice e exige que
vértices conectados por uma aresta possuam cores diferentes.

A condição fundamental é

$$
(u,v)\in E
\quad\Longrightarrow\quad
cor[u]\neq cor[v].
$$

Se essa condição for violada, o grafo não é bipartido. Caso todas as
arestas sejam compatíveis com a 2-coloração, o grafo pode ser dividido
em dois grupos.

A mesma estratégia pode ser implementada utilizando uma matriz de
adjacência ou um TAD baseado em listas de adjacência. A matriz oferece
uma implementação direta e acesso $O(1)$ às arestas, enquanto a lista
de adjacência apresenta melhor complexidade para a BFS em grafos
esparsos:

$$
\boxed{
O(V^2)\quad\text{(matriz)}
}
$$

e

$$
\boxed{
O(V+E)\quad\text{(lista de adjacência)}.
}
$$

\end{document}
